\subsection*{NT-E11. CM Hasse deficits inside a nonabelian Frobenius class}

Let \(E/K\) be an elliptic curve all of whose complex
multiplication is defined over \(K\), and let \(\eta\) be its
normalized unitary CM Hecke character. Thus at good primes
\[
 a_{\mathfrak p}
     =\sqrt{N_K\mathfrak p}
          \{\eta(\mathfrak p)+\overline{\eta(\mathfrak p)}\}.
                                                               \tag{E11.1}
\]
The CM character premise and its precise classical locator
are those of NT-E06 (S-NT-BFK, Proposition 2.5 and the CM
identification in the footnote to Proposition 4.7).
Its nonzero angular infinity type excludes finite order
times a unitary norm power.
Choose any finite Galois \(L/K\), any conjugacy class \(C\),
and the notation \(\delta,Q_C,{\cal L}_C\) of NT-A12.
Include in \(S\) all bad-reduction primes, as well as the
ramification and conductor primes required there.
Assume exactly the cyclic rank-one RH family (A12.3).
Conditional input \(M\) supplies it.

Fix \(c>0\), \(-1/2<\lambda<1/2\), and put
\[
 \begin{aligned}
 \alpha&=\frac14-\frac{\lambda}{2},\qquad
 u_0=\max\{2,c^{1/(1/2-\lambda)}\},\\
 w(u)&=\frac1\pi
       \arccos\left(1-\frac c2 u^{\lambda-1/2}\right)
                  \quad(u\ge u_0).
 \end{aligned}                                             \tag{E11.2}
\]
For either sign, count the primes \(\mathfrak p\notin S\)
with \(u_0<N_K\mathfrak p\le x\), Frobenius class \(C\), and
\[
       0\le 2\sqrt{N_K\mathfrak p}\mp a_{\mathfrak p}
                      <c(N_K\mathfrak p)^\lambda .
\]
Their number \(D_{C;c,\lambda}^{\pm}(x)\) satisfies, for
\(x>\max\{4,u_0\}\),
\[
 \begin{split}
 D_{C;c,\lambda}^{\pm}(x)
   ={}&\delta\int_{u_0}^x\frac{w(u)}{\log u}\,du\\
     &+O_{c,\lambda}\left(
        \delta\sqrt x\log(2x){\cal L}_C(x)
        +\delta x^{5/8+\lambda/4}\sqrt{{\cal L}_C(x)}
                        \right).                         
 \end{split} \tag{E11.3}
\]
All other dependence is explicit in (A12.2). The same formula
holds when only absolute-degree-one primes are counted.
A sufficient uniform relative asymptotic regime is
\[
       {\cal L}_C(x)(\log x)^2=o(x^{1/4+\lambda/2}).
                                                               \tag{E11.4}
\]
For fixed arithmetic data it gives
\[
 D_{C;c,\lambda}^{\pm}(x)
    \sim \frac{\delta\sqrt c}{\pi(3/4+\lambda/2)}
                    \frac{x^{3/4+\lambda/2}}{\log x}.
                                                               \tag{E11.5}
\]

In particular let \(M_C^\pm(x)\) count the primes outside \(S\)
of norm at most \(x\), Frobenius class \(C\), and
\(a_{\mathfrak p}=\pm\lfloor2\sqrt{N_K\mathfrak p}\rfloor\).
For fixed \(E,K,L,C,S\),
\[
 \begin{split}
 M_C^\pm(x)
     &=\frac{\delta}{\pi}
          \int_2^x\frac{u^{-1/4}}{\log u}\,du
                +O_{E,K,L,C,S}\bigl(x^{5/8}\sqrt{\log(2x)}\bigr),\\
 M_C^\pm(x)&\sim
             \frac{4\delta}{3\pi}\frac{x^{3/4}}{\log x}.
 \end{split}                                                   \tag{E11.6}
\]
The integral, rather than only its leading asymptotic term, is
retained in the error-bearing formula. No general non-CM
elliptic \(L\)-function RH is used.

\begin{proof}
Write \(\eta(\mathfrak p)=e^{2\pi i\theta_{\mathfrak p}}\).
For the positive sign the displayed deficit condition says
\[
 \cos(2\pi\theta_{\mathfrak p})
                      >1-\frac c2 q^{\lambda-1/2},
                \qquad q=N_K\mathfrak p .
 \]
It is exactly the arc about zero of normalized length \(w(q)\),
with its outer endpoints omitted. The negative sign gives
the arc of the same length centered at \(1/2\). Equality
at the Hasse endpoint is allowed and is already in these
arcs. All endpoint conventions are covered by NT-A12.

Here is the required uniform control of the moving length.
Set \(v=cu^{-2\alpha}\), so \(0<v\le1\). Direct integration
and differentiation give
\[
 w(u)=\frac1\pi\int_0^v\frac{dz}{\sqrt{4z-z^2}},
 \qquad
 -u w'(u)=\frac{\alpha\sqrt v}{\pi\sqrt{1-v/4}},
 \qquad w(u)\ge\frac{\sqrt v}{\pi}.
 \]
Consequently
\[
 w(u)\asymp_{c,\lambda}u^{-\alpha},\qquad
    0\le-\frac{u w'(u)}{w(u)}
                    \le\frac{2\alpha}{\sqrt3}.            \tag{E11.7}
\]
On a multiplicative interval of ratio at most
\(1+\varepsilon\), with \(0<\varepsilon\le1/2\),
the endpoint lengths therefore have ratio
\(1+O_\lambda(\varepsilon)\).

Partition \((u_0,x]\) by such intervals. The arcs centered
at zero, or at \(1/2\), are nested. Their two endpoint
counts, using (A12.4), bracket the moving count. As in the
complete proof of NT-A13, the total error is
\[
 \begin{gathered}
 O_{c,\lambda}(A/\varepsilon+B\varepsilon),\\
 A=\delta\sqrt x\log(2x){\cal L}_C(x),\qquad
 B=\frac{\delta x^{1-\alpha}}{\log x}.
 \end{gathered}                                           \tag{E11.8}
\]
For completeness the first term follows from
\(\sum_j\sqrt{u_j}\ll\varepsilon^{-1}\sqrt x\).
For the second, (E11.7) bounds the difference of the
upper and lower integral sums by
\[
 O_\lambda\left(\delta\varepsilon
                    \int_{u_0}^x\frac{w(u)}{\log u}\,du\right)
             \ll_{c,\lambda} B\varepsilon .
 \]
Their common limiting integral is precisely the first
term in (E11.3), so no small-angle expansion has been
substituted into this comparison. The initial bounded
range is uniform after absorption by
\([K:\mathbf Q]\le\delta n_L\) and
\(\log Q_C\ge n_L\log4\).

Take \(\varepsilon=\min\{1/2,\sqrt{A/B}\}\).
Then (E11.8) is \(O_{c,\lambda}(A+\sqrt{AB})\), and
\[
 \sqrt{AB}\ll
          \delta x^{3/4-\alpha/2}\sqrt{{\cal L}_C(x)}
       =\delta x^{5/8+\lambda/4}\sqrt{{\cal L}_C(x)}.
 \]
This proves (E11.3). Using the degree-one form of (A12.4)
throughout proves its degree-one version.
Since \(w(u)\sim(\sqrt c/\pi)u^{-\alpha}\), the integral
has asymptotic
\[
 \frac{\delta\sqrt c}{\pi(1-\alpha)}
                          \frac{x^{1-\alpha}}{\log x}.
 \]
The relative errors are bounded by a constant times
\(R(x)+\sqrt{R(x)}\), where
\(R(x)={\cal L}_C(x)(\log x)^2x^{\alpha-1/2}\).
This proves (E11.4)--(E11.5).

For (E11.6), take \(c=1\), \(\lambda=0\). For any integer
\(a\le y\), \(a=\lfloor y\rfloor\) is equivalent to
\(a>y-1\); this includes integral \(y\).
Apply this with \(a=a_{\mathfrak p}\), or
\(-a_{\mathfrak p}\), and \(y=2\sqrt q\).
The floor event is exactly the strict deficit condition
in (E11.3), even when \(q\) is a square.
The primes of norm exactly two contribute at most
\([K:\mathbf Q]\le\delta n_L\), absorbed by its error.
The elementary expansion on \(u\ge2\) is
\[
 w(u)=\frac1\pi u^{-1/4}+O(u^{-3/4}).
 \]
Its integral remainder is \(O(x^{1/4})\) and is absorbed
by the first error in (E11.3).
For fixed arithmetic data
\({\cal L}_C(x)=O_{E,K,L,C,S}(\log(2x))\).
Thus the second error is \(O(x^{5/8}\sqrt{\log(2x)})\);
the first, \(O(\sqrt x\log^2(2x))\), is smaller for
sufficiently large \(x\), and is absorbed after adjusting
the fixed-data constant. Finally
\(\int_2^x u^{-1/4}du/\log u
 \sim(4/3)x^{3/4}/\log x\), which proves (E11.6).
\end{proof}

\paragraph{Scope and geometric placement.}
The finite étale cover may have an arbitrary nonabelian
Galois group; for example it may be a division-field
cover attached to another, non-CM elliptic curve.
The trace in this theorem is still the trace of the
specified CM curve \(E\). The proof couples the finite
Frobenius condition to its rank-one CM torus by the cyclic
intermediate covers of NT-A12. It does not infer GRH for
the standard \(L\)-function of that second curve, or a
Koblitz prime-order statement. The assumption that all
CM is defined over \(K\) is retained; no unproved descent
factor is inserted into this simultaneous formulation.
At \(\lambda=-1/2\), the positive exponent in (E11.4)
vanishes and the argument gives no asymptotic; the
finite norm-form gate in NT-E10 remains applicable.
For general \(\lambda\), the exact arcsine integral in
(E11.3) is essential to the stated error, as explained
in NT-E09. This is a newly proved formulation, with no
claim of literature priority.
Work/run: \texttt{GRH\_NIGHT/number\_theory},
\texttt{cm\_frobenius\_deficits.tex}, NT-E11, also
\texttt{section.tex}; full analytic input and its
classical locators are (A12.3), NT-A08 and NT-E06.
